\chapter{Program, Scope, and Standing Conventions}
\label{ch:introduction}

This manuscript asks what collective inference is available to agents that are local sections of a
belief and a model associated bundle of one principal bundle, placed at one base point and interacting
through independent links under one normalized joint law, and which parts of that inference survive
exact coarse-graining of the agents into meta-agents. The answer has four main parts. The whole
population has one evidence bound, an identity in the extended reals whose gap vanishes exactly at the
posterior as a measure (\Cref{thm:elbo-extended-gap,cor:elbo-bound-tightness}). Each agent's local
objective is a conditional free energy whose differential is that of the collective free energy
(\Cref{thm:obs-local-global-potential}). Coarse channels of the agent network preserve the evidence and
decompose the free energy exactly (\Cref{thm:rg-exact-coarse-vfe}), with an effective action for the
retained coordinates (\Cref{thm:rg-effective-action}). A finite-depth family forms successive
meta-agents under free-energy descent with learning, in a stated scope
(\Cref{thm:finite-depth-adaptive-formation}). \Cref{sec:intro-reading-map} maps these results to the
chapters, and \Cref{sec:intro-conventions} fixes, once for the whole manuscript, the defaults and the
claims it declines.

\section{Program aim}
\label{sec:intro-program-aim}

The program is a participatory multi-agent theory. Agents are local sections of associated bundles,
collectives form through the agents' coupled inference, and parents act back on their children through
normalized laws rather than through quantities that read the children's beliefs. Three requirements
follow from that aim, and the manuscript is organized around them. Every collective object must be a
complete probability object, constructed from the agents' mechanisms and interaction records or selected
explicitly, never an additional agent and never a sum of local objectives that happens to look like one.
Every statement about a parent must hold for the normalized parent law that coarse-graining actually
produces. And formation of a parent must be a theorem about a declared family, with its scope stated,
rather than a reading of a formula.

The present phase fixes one base point $c_*$ and settles the theory there before any transport across the
base is considered. At $c_*$ the exact normalized parent law exists for every coarse channel of the
agent network (\Cref{thm:rg-pointwise-parent-datum}), and \Cref{thm:finite-depth-adaptive-formation} gives
a formation witness under a supplied hierarchy of interaction strengths. On a homogeneous four-cycle a
random initial source assignment already makes a collective form from every Gaussian initial law, with
source learning idle (\Cref{prop:phil-seeded-contrast-formation}). At each sufficiently small fixed
temperature, slow source learning selects one matching in the assignment law for almost every Gaussian
initial mean, while the frozen-prior pairing difference vanishes at long times
(\Cref{conj:phil-contrast-formation}). How a
parent's own dynamics are determined by the network it summarizes is
\Cref{open:rg-pointwise-parent-dynamics}, the defining open problem of the phase.

\section{Objects and the local-first architecture}
\label{sec:intro-objects}

Agent $i$ is the object $\mathcal A^i=(\mathcal C_i;q_i^b,q_i^m,u_i)$ of \Cref{def:geo-agent}: a belief
section and a model section of the associated bundles $\mathcal E_b$ and $\mathcal E_m$ of one principal
$G$-bundle, expressed in one common frame $u_i$, a gauge choice by \Cref{prop:geo-common-frame-gauge}. The
two channels may carry inequivalent representations, and cross morphisms between them are additional
structure (\Cref{def:geo-cross-morphisms}). Interactions are normalized record kernels $K_a$ on the scopes
of an interaction multigraph, and the links $\Theta_e^x$ on its edge copies are independent data with loop
holonomy and per-edge invariant classes (\Cref{prop:geo-edge-twist-invariance}). The primitive descriptive
data are the agents, their local recognition laws, their evaluated mechanisms, and the record kernels.
Population laws enter only by construction (\Cref{prop:gen-exact-normalization}) or by an explicit
selection (\Cref{hyp:prob-measurable-coupling-selector}). The architecture is
\begin{equation}
\left\{
\mathcal A^i,G_{\theta,i},\mathbb Q_{i,o,X},K_a
\right\}_{i\in V,a\in\mathfrak A}
\longrightarrow
\begin{cases}
\mathbb P_{\theta,V} & \text{normalized construction},\\
\mathbb Q_{V,o,X}\in\operatorname{Cpl}(\mathbb Q_{\bullet,o,X})
& \text{measurable selection},
\end{cases}
\longrightarrow
\Lelbo_V,
\label{eq:intro-local-to-global-architecture}
\end{equation}
so the constructed joint and the selected coupling precede every local and collective objective. On the
evidence domain of \Cref{hyp:elbo-evidence-domain} there is one identity,
\begin{equation}
\log p_{\theta,V}(o\given X)
=\Lelbo_V(\mathbb Q_{V,o,X};X,o)
+\KL\bigl(\mathbb Q_{V,o,X}\big\Vert\boldsymbol\Pi_{\theta,V,o,X}\bigr),
\label{eq:elbo-intro-identity}
\end{equation}
with equality of the bound and the log evidence exactly when the recognition coupling is the posterior as
a measure, which matching coordinate marginals does not achieve (\Cref{cor:elbo-bound-tightness}). The
identity is invariant under passive re-expression of the frames when the kernels are covariant
(\Cref{prop:gen-product-evidence-invariance,cor:elbo-independent-frame-invariance}). Excluding the
posterior from a recognition family makes every pointwise gap positive, but it does not make the optimized
gap positive or attained (\Cref{sec:restrict-principle}).

\section{Results and reading map}
\label{sec:intro-reading-map}

Part I builds the single-point foundations. \Cref{ch:geometry} fixes the geometry at $c_*$: the active gauge
group is one copy of $G$, links carry per-edge twist classes, and three notions of agreement are separated.
\Cref{ch:probability} fixes kernels before densities, proves the chain rule used everywhere, and
distinguishes sample, law-point, and configuration levels. \Cref{ch:generative} builds the normalized joint
on arbitrary record cycles and states the typing prohibition. \Cref{ch:elbo} derives the evidence bound and
the local-collective potential identity, with observations entering as interaction records
(\Cref{thm:obs-agent-interaction-equivalence}). \Cref{ch:selective-communication} places attention inside
the fixed joint as a latent source label (\Cref{prop:obs-attention-elbo}), builds selective communication,
and records its correspondence with transformer attention (\Cref{tab:obs-transformer-correspondence}).
\Cref{ch:relational-inference} develops fiber information geometry, Fisher contraction under Markov
morphisms (\Cref{thm:cg-fisher-contraction}), and inference histories measured by a Fisher duration that is
invariant under reparameterization (\Cref{thm:hist-fisher-clock-invariance}). \Cref{ch:two-timescales} types
fast beliefs and slow models (\Cref{def:obs-fast-slow-blocks}) and proves the reduced slow flow
(\Cref{prop:obs-rate-separated-flow}).

Part II coarse-grains at one base point. \Cref{ch:coarsegraining} separates Markov coarse channels, energy
restriction, and parameter aggregation (\Cref{sec:cg-three-operations}) and proves the exact coarse free
energy and the pointwise parent datum (\Cref{thm:rg-exact-coarse-vfe,thm:rg-pointwise-parent-datum}).
\Cref{ch:network-state} gives the gauge-retained network state, the kernel--holonomy fixed sectors
(\Cref{prop:cg-kernel-holonomy}), and the holonomy-conditioned marginal law
(\Cref{thm:cg-holonomy-kl-marginal}). \Cref{ch:agent-network-rg} is exact coarse-graining of the agent
network: effective action, residual, finite-step multipliers, and the native growth bound
(\Cref{thm:rg-effective-action,thm:rg-bounded-action-calculus,prop:rg-interaction-native-growth-bound}),
with the retained state collected in \Cref{tab:rg-retained-state} and a parent-sensitive mechanism with
path laws and an exact finite-memory cost (\Cref{sec:rg-memory-parent-mechanism}).
\Cref{ch:full-graph-meta-agent-vfe} assembles the multiscale free energy and closes a tower of priors
(\Cref{sec:full-graph-closing-tower}).

Part III turns to formation and interpretation. \Cref{ch:finite-depth-formation} constructs, for every
finite depth, a family in which successive coherence episodes come from linear relaxation of a prepared mean
through a supplied hierarchy of interaction strengths; source jumps and learning are perturbations the proof
controls, not the cause of formation, and a separate witness produces a four-cycle endogenously
(\Cref{prop:rg-formation-endogenous-cycle}). \Cref{ch:philosophy} states what the mathematics does and does
not support, and \Cref{ch:related-work} positions the construction. The appendices hold measure-theoretic
lemmas (\Cref{app:measure-lemmas}), the formation proofs (\Cref{app:finite-depth-formation-proofs}), the
notation (\Cref{app:notation-contract}), the open problems (\Cref{app:claim-ledger}), and the numerical index
(\Cref{app:numerical-provenance}).

Two supplements accompany the manuscript, and no result here depends on them. Supplement G develops the
multivariate-Gaussian realization. Supplement B collects base transport, covers and connection descent,
multi-point designs, infinite networks with their asymptotic apparatus, and the lattice cluster.

A first reading can take \Cref{ch:geometry,ch:probability,ch:generative} for the objects and the joint, then
\Cref{ch:elbo} for the bound and \Cref{ch:selective-communication} for attention, then
\Cref{ch:coarsegraining,ch:agent-network-rg} for exact coarse-graining and
\Cref{ch:finite-depth-formation} for formation; proofs, witnesses, and the ledger can be recovered through
their cross-references on a second reading.

\section{Standing conventions and scope}
\label{sec:intro-conventions}

Everything takes place at one base point. The design is $\{c_*\}$, design and incidence indices are dropped,
and every agent's frame, beliefs, and models are points of the fibers over $c_*$. Base maps and base curves,
neighborhoods, connection descent, transport along the base, section spaces over $\mathcal C$, and continuum
law theory are deferred to the base supplement named in \Cref{sec:intro-reading-map}.

The following defaults hold throughout, and a statement mentions only departures from them. Measurable spaces
are nonempty and standard Borel, and $\mathcal P(\mathsf S)$ denotes the probability laws on a measurable
space $\mathsf S$. Laws and kernels are normalized and measurable. A coarse channel is a Markov
kernel on latent coordinates that does not depend on the recognition law, leaves observations unchanged, and
is vertical over the identity of $c_*$. An admitted observation has finite positive evidence. Structural data
are parameters unless a configuration prior or a structural channel is declared, and the declaration names
its law. The interaction multigraph at $c_*$ is a structural parameter whose edge copies come in reversal
pairs; it contains every copy in the support of a receiving row's source prior and, unless a model adds
further copies, is the reversal closure of the union of these supports. Each support is complete unless
masked, and supports need not be symmetric: a receiver may attend to a source that does not attend to it. Gibbs attention weights are strictly positive on positive-prior edges when their
costs are finite; the extended-cost rows of \Cref{prop:full-graph-gibbs-rows} can assign zero weight to
infinite-cost edges. Dynamic structure is carried by attention weights and link values. Discrete creation and deletion of edges is not part of the model; it is
\Cref{open:obs-edge-creation}. Coarse multigraphs are induced structure. Absolute continuity of a recognition
law with respect to the posterior is not a default: it is clause (H3) of \Cref{hyp:elbo-evidence-domain}, and
its failure is the singular branch of the extended bound.

Three transports must remain distinct. A pointwise comparison $T_{ij}$ relates two frames in the same
principal fiber and involves no curve (\Cref{prop:geo-frame-comparisons}). A link $\Theta_e^x$ is an
independent variable on an edge copy of the interaction multigraph (\Cref{def:geo-graph-links}). Parallel
transport along a base curve is not used. Graph links are not base transport, and graph-link holonomy is
neither the holonomy of a base connection nor a property of the topology of the principal bundle.

The word coarse has two uses that must not be merged. A Markov channel discards sample information while
preserving normalization; a hard restriction substitutes a lower-dimensional configuration into an energy.
They can agree in a special realization, but a shared word transfers no theorem, and parameter aggregation
inside a closed family is a third operation (\Cref{sec:cg-three-operations}). Likewise, an inference history
is an oriented unparameterized orbit in a configuration space, an ordered interaction-record path retained by
exact coarse-graining is a sample coordinate, coarse-graining depth is a scale index, and a Fisher duration is
a statistical length.

\paragraph{Scope.}
The manuscript declines the following claims. None of the inference index, flow parameter, transition index,
ordered path coordinate, coarse-graining depth, or Fisher duration is identified with physical time; such an
identification would need a named target clock, a calibration map with an uncertainty model, causal and
compositional consistency, and discriminating evidence against alternative parameterizations. The contextual
base is not spacetime and carries no Lorentzian structure. Sections, laws, and meta-agents are mathematical
objects: calling a boundary node an agent is a naming decision, no claim of autonomous agency or of what
exists is made, and population laws, population samples, and blocks are not agents. No canonical coupling
selector, partition, or membership selector is asserted. No renormalization semigroup with rescaling, fixed
points, or beta functions is constructed; Supplement B holds that asymptotic apparatus. The identification of
any limiting object with a physical law is \Cref{claim:physical-law-identification}. Declining a claim is not
refuting it, and a refuted statement appears instead as an established negative result. \status{NOT-CLAIMED}

\section{Status of statements}
\label{sec:intro-status}

Every definition, standing hypothesis, lemma, proposition, theorem, corollary, conjecture, open problem, and
claim-bearing remark is a numbered statement. Definitions, standing hypotheses, conjectures, and open problems
carry their status in their heading type. Every lemma, proposition, theorem, corollary, and claim-bearing
remark carries exactly one of three tags, one for all of its conclusions. Text that motivates, summarizes, or
points carries no tag and makes no new claim. The fourth tag of \Cref{tab:status-taxonomy}, for declined
claims, appears once in each scope paragraph and nowhere else. Open problems are collected in \Cref{tab:open-problems}, and
numerical statements with their checks in \Cref{tab:numerical-index}.

\begin{sciencetable}{Status of statements.}
\label{tab:status-taxonomy}
\begin{tabularx}{\textwidth}{>{\raggedright\arraybackslash}p{0.24\textwidth}>{\raggedright\arraybackslash}X}
\toprule
\textbf{Tag or heading} & \textbf{Meaning} \\
\midrule
\status{ESTABLISHED} & Proved in this manuscript. A refutation is an established negative statement. \\
\status{STANDARD} & A known result, cited to a checked primary source with a theorem or section locator. \\
\status{NUMERICAL} & Supported only by a named computation, indexed with its control and seed; computation is not proof. \\
Definition & A declared type, construction, or convention; nothing is proved. \\
Hypothesis & A standing restriction adopted for a chapter or part; it states what it excludes. \\
Conjecture & A precise unproved statement with a named observation that would refute it. \\
Open problem & An unsettled question with what holds, what is unproved, and what would settle it. \\
\status{NOT-CLAIMED} & A claim declined in a scope paragraph; declining is not refuting. \\
\bottomrule
\end{tabularx}
\end{sciencetable}

\section*{Physicist's summary}

The manuscript studies agents that are local sections of belief and model bundles, all at one base point,
coupled through independent links and interaction records under a single normalized joint law. There is one
evidence bound for the whole population, the local objectives are exact conditional pieces of it, and
coarse-graining the agent network into meta-agents is done with Markov channels that preserve the evidence
and split the free energy exactly. Formation of meta-agents is proved for a declared finite-depth family in a
stated regime and conjectured more generally. Physical time, spacetime, agency, and transport across the base
are not claimed; they are either deferred to the base supplement or declined, and the defaults stated above
are not repeated in later chapters.
